% v21: exact sector control for positive two-high branches with L=3 or 4.
\section{小さい低辺係数を持つ正枝の符号と分母順序}
\label{ps21:section}

Book の低三角形では，正枝の非零低 entry は $-3$ または $-4$ である．
この二つの値に限ると，pair equation の係数の符号は高分母の順序と結び付き，
低 $D$ が負となる側には reverse parameter の強い上界がある．
以下の局所補題には分母の奇数性，node 条件，clutching は不要である．

\begin{lemma}[正枝の二つの sector]\label{ps21:positive-sector}
正の整数二高分母標準形において
\[
 \delta=T-S\in\{2,3\},\quad d=\delta-1,\quad
 h=S-r,\quad k=T-r-1=h+d,
\]
\[
 \alpha=\frac{sh}{r},\quad\beta=\frac{zk}{r+1},\quad
 F=(2r+1)h-r^2d
\]
と置く．このとき $F\ne0$ であり，次が成り立つ．
\begin{enumerate}
\item $F>0$ なら $s\le z$，$D_D<0<D_C$，$a<b$ である．
\item $F<0$ なら $D_D\ge0$ であり，
\begin{equation}\label{ps21:h-bound}
 h\le
 \begin{cases}
 \lfloor(r-1)/2\rfloor,&\delta=2,\\
 r-1,&\delta=3.
 \end{cases}
\end{equation}
また $a>b$ であるか，次の二つの例外のいずれかである：
\begin{equation}\label{ps21:equal-base-points}
 (r,s,z,S,T)=(2,2,1,3,6),\qquad (3,3,2,4,6).
\end{equation}
この二点ではそれぞれ $a=b=7,13$，$D_D=0$，$D_C=2,1$ である．
\end{enumerate}
特に $D_C\le0$ なら $F<0$，$a>b$ であり，
\begin{equation}\label{ps21:T-bound}
 T\le
 \begin{cases}
 \lfloor(3r+3)/2\rfloor,&\delta=2,\\
 2r+2,&\delta=3.
 \end{cases}
\end{equation}
\end{lemma}

\begin{proof}
正枝の既知の不等式 $rT>(r+1)S$ により $0<h<rd$ である．
補題~\ref{ml12:dichotomy} と式~\eqref{td12:bezout} より，
$\alpha,\beta$ は正の整数であり，
\begin{equation}\label{ps21:two-identities}
 F\alpha\beta=(\beta-\alpha-1)hk,\qquad
 M\beta-N\alpha=1,
 \quad M=rs+1,
 \quad N=(r+1)z+1.
\end{equation}
特に $\gcd(\alpha,\beta)=1$ である．$F=0$ なら
$2r+1\mid r^2d$ となるが，$\gcd(2r+1,r^2)=1$ と $d\in\{1,2\}$ に反する．

まず $F>0$ とする．式~\eqref{ps21:two-identities} により
$j=\beta-\alpha\ge2$ である．$s\ge z+1$ と仮定すると，
第二式から
\[
 M\{r\beta-(r+1)\alpha\}
 =r+\{r(r+1)(z-s)-1\}\alpha<0
\]
を得る．従って $\alpha>rj$，$\beta>(r+1)j$ となる．
一方，$h<rd$，$k<d(r+1)$，$F\ge1$ なので
\[
 j-1=\frac{F\alpha\beta}{hk}
 >\frac{j^2}{d^2}\ge\frac{j^2}{4}\ge j-1,
\]
ここで最後の不等式は $(j-2)^2\ge0$ である．これは矛盾であり，$s\le z$ である．
従って $D_D=s-z-1<0$，また
$T/(r+1)>S/r$ から $D_C=\delta-sS/r+zT/(r+1)>0$ となる．
さらに
\[
 a-b=rs-(r+1)z+r\alpha-(r+1)\beta<0
\]
である．

次に $F<0$ とする．整数 $h$ に対する
$h<r^2d/(2r+1)$ は式~\eqref{ps21:h-bound} を与える．
$j=\alpha-\beta\ge0$ と置くと，
\[
 E=-\frac{F}{hk}>0,\qquad E\alpha\beta=j+1.
\]
また
\[
 -F-k=(r+1)\{d(r-1)-2h\}\ge0
\]
なので $E\ge1/h$ であり，
\[
 E\beta=\frac{j+1}{\beta+j}\le1
 \quad\Longrightarrow\quad \beta\le h<r.
\]
$rk-(r+1)h=rd-h>0$ と $\alpha\ge\beta$ により
\[
 \frac sz=\frac\alpha\beta\frac{rk}{(r+1)h}>1,
\]
従って $D_D\ge0$ である．$j\ge1$ なら
$M\beta-N\alpha=1$ より $M>N$，すなわち
$rs-(r+1)z\ge1$ である．ゆえに
\[
 a-b=rs-(r+1)z+rj-\beta>0.
\]

残る $j=0$ では $\alpha=\beta=1$ である．
$rs-(r+1)z=1$ なので，$m=s-z\ge1$ と置くと
$z=rm-1$，$s=(r+1)m-1$ である．
$\alpha=sh/r=1$ から $s\le r$ となり，$m=1$，$s=r$，$z=r-1$ を得る．
従って $h=1$，$k=(r+1)/(r-1)$ であり，
$r-1\mid2$ から式~\eqref{ps21:equal-base-points} の二点に限られる．
直接代入によって表示した分母と低 $D$ を得る．
最後に $D_C\le0$ では $F>0$ とこの二点が排除され，
式~\eqref{ps21:h-bound} に $T=r+h+\delta$ を代入すると
式~\eqref{ps21:T-bound} を得る．
\end{proof}

\begin{corollary}[Book の $D$ 周期に沿う高分母の順序]
\label{ps21:book-denominator-order}
定理~\ref{bd20:nonzero-cycle} の book 低三角形を考える．
二つの高頂点を $0,1$ とし，低 row の正 entry が高 ray $0$ を向く側を
next とする有向 node 三角形を用いる．この向きに沿う各低 $D$ が正なら
\[
 q_{i0}<q_{i1}\qquad(i=2,3,4),
\]
全て負なら三つの不等号は全て逆になる．
\end{corollary}

\begin{proof}
正枝では，非零低 $D$ が一正一負であることと補題~\ref{ps21:positive-sector} から
\[
 D_C<0\Longleftrightarrow F<0\Longleftrightarrow a>b,
 \qquad
 D_C>0\Longleftrightarrow F>0\Longleftrightarrow a<b
\]
となる．負枝でも式~\eqref{bd20:negative-low-D} と
定理~\ref{ml12:denominator-descent} の証明により，同じ $D_C$ と $a-b$ の対応がある．
canonical $A=0$ なら next は $C$，$a=q_{i0}$，$b=q_{i1}$ である．
canonical $A=1$ なら next は $D$，$a=q_{i1}$，$b=q_{i0}$ であるため，
$D$ の符号と分母順序が同時に反転する．
対向 $D$ の相殺と各 core の一正一負性から，node 三角形上の next の $D$ は
三つ全て同符号である．これで両方の場合が従う．
\end{proof}

この系は中心二 core に接続する三つずつの分母を比較する必要条件である．
中心が四高分母を持つことを用いる追加の議論なしに，book 全体の不存在を結論しない．
